\documentclass[a4paper,10pt]{article}
\usepackage[spanish]{babel}
\usepackage{
			ucs
			,amsfonts
			,amssymb
			,acronym
			,verbatim
			,rotating
			,graphicx
			,enumerate
			,fancyhdr
            ,mathrsfs
            ,array
			}

\usepackage{tikz}
\usetikzlibrary{svg.path}

\pagestyle{fancy}
\addtolength{\textheight}{1.0in}
\addtolength{\textwidth}{0.5in}
\addtolength{\voffset}{-1.0in}

\newcommand{\cat}[1][]{%
    \tikz \fill [scale=1ex/30,yscale=-1,#1] svg "M495.959,867.638c-3.318-4-1.121-7.686,7.098-11.904c7.322-3.76,8.223-4.871,11.365-14.049
		c5.646-16.496,9.951-21.084,24.289-25.893c4.883-1.637,7.111-3.217,6.732-4.771c-3.148-12.889-1.992-16.166,4.48-12.701
		c2.174,1.162,3.924,0.754,6.943-1.619c4.547-3.576,5.105-3.137,6.68,5.264c0.725,3.859,0.445,6.785-0.791,8.273
		c-1.186,1.43-1.496,4.199-0.816,7.295c0.895,4.07,0.186,6.26-3.836,11.861c-6.01,8.367-6.746,18.059-1.752,23.053
		c3.975,3.973,4.066,6.363,0.248,6.363c-3.869,0-8.201-4.926-8.201-9.322c0-4.723-4.635-6.33-7.391-2.564
		c-2.844,3.889-2.738,7.115,0.23,7.115c1.312,0,2.387,1.074,2.387,2.387c0,1.92-2.465,2.385-12.635,2.385
		c-6.947,0-13.047-0.668-13.551-1.484c-0.541-0.875-3.803-0.502-7.951,0.912c-7.744,2.639-11.631,5.996-8.896,7.686
		c0.916,0.566,1.668,1.703,1.668,2.527C502.26,870.757,498.104,870.222,495.959,867.638z";~%
}


\fancyhf{}
\fancyhead[L]{{\bfseries Algoritmos y Estructuras de Datos. Pr\'actica 2.}}
\fancyfoot[LE,RO]{\thepage}

\fancypagestyle{firstpage}{
\fancyhf{}

\fancyhead[L]{\includegraphics[width=\textwidth]{header.png} 
}
\fancyfoot[LE,RO]{\thepage}
}

\thispagestyle{firstpage}

\begin{document}

~\\ \\ \\ \\
\begin{itemize}

\item[3]
    \begin{itemize}
        \item[d-] \begin{verbatim} Gamma := tar;\end{verbatim} INV\'ALIDA. No se puede asignar un valor a una constante.
        \item[e-] \begin{verbatim} rango := respuesta;\end{verbatim} INV\'ALIDA. Los tipos de las variables no concuerdan.
        \item[g-] \begin{verbatim} tar := mx9;\end{verbatim} V\'ALIDA. El valor de \verb=mx9= se convierte autom\'aticamente a real (al reves no se puede).  
        \item[h-] \begin{verbatim} simbolo := `t`;\end{verbatim} V\'ALIDA. Los tipos son correctos y se puede realizar la asignaci\'on.
    \end{itemize}

\item[4.b-] ~~
    \begin{table}[!h]
      \begin{tabular}{|c||c|c|c|}
        \hline
        Nro. de Linea & \verb=x= & \verb=y= & \verb=z= \\
        \hline \hline
        1 & 27 & -- & -- \\
        2 & 27 & 343 & -- \\
        3 & 27 & 343 & 70 \\
        4 & 30 & 343 & 70 \\
        5 & 30 & 100 & 70 \\
        6 & 200 & 100 & 70 \\
        7 & 200 & 200 & 70 \\
        \hline
      \end{tabular}
    \end{table}

\item[5.c-]
\begin{small}
\begin{verbatim}
for i in range(0,5) :
    for j in range(0,  4 - (i \% 4)) :
        print(' ', end="")
    for j in range(0, (2 * (i \% 4)) + 1) :
        print('*', end="")
    for j  in range(0, 4 - (i \% 4)) :
        print(' ', end="")
    print('')
\end{verbatim}
\end{small}

\item[7]
\begin{small}
\begin{verbatim}
total = int(input('Ingrese la cant. de dinero: \$ '))
print('Billtes de \$100: ', total // 100)
total = total \% 100
print('Billtes de \$50: ', total // 50)
total = total \% 50
print('Billtes de \$20: ', total // 20)
total = total \% 20
print('Billtes de \$10: ', total // 10)
total = total \% 10
print('Billtes de \$5: ', total // 5)
total = total \% 5
print('Billtes de \$2: ', total // 2)
total = total \% 2
print('Monedas de \$1: ', total)
\end{verbatim}
\end{small}

\item[8]
\begin{small}
\begin{verbatim}
segundos = int(input('Ingrese la cant. de segundos: '))
h, m, s = (segundos // 60) // 60, (segundos // 60) % 60, ((segundos % 60) % 60)
print(h, m, s)  
\end{verbatim}
\end{small}

\end{itemize}
\end{document}
